% answers_short.tex —— 2009 年·普通高等学校招生考试·湖南卷·文科 · 简短答案速查
% 由 tools/make_short_answers.py 自 answers.tex 的 \gksolhead/\gkanswer 提取，
% 并与 2009/2009_数学真题_湖南文_答案要点.md 逐题比对后生成，勿手改（改源 MD 后重跑）。

\begin{gknotice}
\textbf{使用说明}：本节只给最终答案便于快速核对，完整推导见解析版。答案为编写者独立推导，\textbf{非官方标准答案}。
\end{gknotice}

\gkpart{选择题}
\begin{center}\begin{tabular}{ccccccccc}
\toprule
题号 & 1 & 2 & 3 & 4 & 5 & 6 & 7 & 8 \\
\midrule
答案 & D & B & C & A & B & C & A & C \\
\bottomrule
\end{tabular}\end{center}

\gkpart{填空题}
\begin{itemize}[leftmargin=2.6em]
  \item[\textbf{9.}] $12$。
  \item[\textbf{10.}] $2\sqrt2$。
  \item[\textbf{11.}] $6$。
  \item[\textbf{12.}] $120$。
  \item[\textbf{13.}] $2$。
  \item[\textbf{14.}] $2$；$(\sqrt2,\sqrt3)$。
  \item[\textbf{15.}] $x=1+\dfrac{\sqrt3}{2}$，$y=\dfrac{\sqrt3}{2}$。
\end{itemize}

\gkpart{解答题}
\begin{itemize}[leftmargin=2.6em]
  \item[\textbf{16.}] （1）$\dfrac14$；（2）$\theta=\dfrac\pi2$ 或 $\dfrac{3\pi}{4}$。
  \item[\textbf{17.}] （1）$\dfrac16$；（2）$\dfrac{19}{27}$。
  \item[\textbf{18.}] （1）证明见解析；（2）$\dfrac{\sqrt{21}}{8}$。
  \item[\textbf{19.}] （1）$b=-6$；（2）定义域 $(2,+\infty)$，值域 $(-\infty,8)$。
  \item[\textbf{20.}] （1）$\dfrac{x^2}{8}+\dfrac{y^2}{4}=1$；（2）$\left[-\dfrac{\sqrt3-1}{2},\dfrac{\sqrt3-1}{2}\right]$。
  \item[\textbf{21.}] （1）是 $B-$数列；（2）②$\Rightarrow$④ 为真命题；（3）证明见解析。
\end{itemize}
